\section{A resonance theorem uniform in the integer order}
\label{sec:resonance}

The Lerch expansion has two singular summands when the order approaches a
positive integer. Their sum is regular. Classical formulas describe this
cancellation for a fixed integer order; the useful additional point here is a
normalization that controls \emph{every} integer order at once. It allows the
order, its distance from an integer, and the distance of the argument from
one to vary simultaneously.

\subsection{The classical identity and the normalized remainder}
For $a>0$ and $t>0$ put
\begin{equation}\label{res:F}
 F(s,t,a)=e^{-at}\Phi(e^{-t},s,a)
 =\sum_{n=0}^{\infty}\frac{e^{-(n+a)t}}{(n+a)^s}.
\end{equation}
In particular $F(s,t,1)=\Li_s(e^{-t})$. The classical Lerch expansion is
\begin{equation}\label{res:classical}
 F(s,t,a)=\Gamma(1-s)t^{s-1}
       +\sum_{j=0}^{\infty}\zeta(s-j,a)\frac{(-t)^j}{j!},
 \qquad 0<t<2\pi,
\end{equation}
initially for $s\notin\{1,2,\ldots\}$. We use the locally normally convergent
form of this identity; see DLMF~25.14.3\_1 \cite{dlmflerch}. Integer-order
polylogarithm derivatives obtained by pairing the poles are already treated
by Bailey and Borwein \cite{baileyborwein}. The prior boundary continuation
in ProveIt also uses this cancellation to establish sharp endpoint regularity
\cite{repoBoundary}. We do not claim those facts as new.

Let $k\geq0$ be an integer and write $s=k+1+\varepsilon$. Empty sums and
products have their usual values. Define
\begin{equation}\label{res:normalized}
 \mathcal R_k(\varepsilon,t,a)=
 \frac{k!}{(-t)^k}\left[
 F(k+1+\varepsilon,t,a)
 -\sum_{j=0}^{k-1}\zeta(k+1+\varepsilon-j,a)\frac{(-t)^j}{j!}
 \right].
\end{equation}
This subtraction removes the regular terms that precede the resonance. It
is an exact definition, even when numerical evaluation by direct subtraction
would lose many digits. Introduce
\begin{align}
 H_k^{(r)}&=\sum_{j=1}^k j^{-r},\qquad H_k=H_k^{(1)},\notag\\
 T&=\log(1/t)+H_k-\gamma,\label{res:T}\\
 G_a(\varepsilon)&=\zeta(1+\varepsilon,a)-\frac1\varepsilon
  =\sum_{r=0}^{\infty}\frac{(-1)^r\gamma_r(a)}{r!}\varepsilon^r.
 \label{res:G}
\end{align}
The apparent singularity of $G_a$ is removable. Throughout the theorem below
$T$ is positive and tends to infinity. No separate upper bound on $k$ is
imposed.

\begin{definition}[Centered resonance coefficients]\label{res:coeffdef}
For $|w|<1$, define
\begin{align}
 E_k(w)&=\exp\left\{\sum_{r=2}^{\infty}
       \frac{\zeta(r)+(-1)^rH_k^{(r)}}{r}w^r\right\}
       =\sum_{j=0}^{\infty}b_{k,j}w^j.\label{res:E}
\end{align}
Thus $b_{k,0}=1$ and $b_{k,1}=0$. If
$c_{k,r}=\zeta(r)+(-1)^rH_k^{(r)}$, the coefficients can be generated by
\begin{equation}\label{res:recurrence}
 j b_{k,j}=\sum_{r=2}^j c_{k,r}b_{k,j-r}\qquad(j\geq2).
\end{equation}
\end{definition}

These coefficients are polygamma polynomials. More explicitly,
\begin{equation}\label{res:polygamma}
 c_{k,r}=(1+(-1)^r)\zeta(r)
       -\frac{\psi^{(r-1)}(k+1)}{(r-1)!},\qquad r\geq2.
\end{equation}
Equation~\eqref{res:polygamma} follows directly from the polygamma series.
It also gives the complete Bell-polynomial description of $b_{k,j}$, with
cumulants $(r-1)!c_{k,r}$ and zero first cumulant.

\begin{lemma}[Exact cancellation]\label{res:exact}
For $|\varepsilon|<1/2$ and $0<t<2\pi$,
\begin{align}
 \mathcal R_k(\varepsilon,t,a)
  &=\frac{1-e^{-T\varepsilon}E_k(\varepsilon)}{\varepsilon}
          +G_a(\varepsilon)+B_k(\varepsilon,t,a),\label{res:pair}\\
 B_k(\varepsilon,t,a)
  &=\sum_{\ell=1}^{\infty}
     \frac{k!(-t)^\ell}{(k+\ell)!}
     \zeta(1+\varepsilon-\ell,a).\label{res:B}
\end{align}
The quotient at $\varepsilon=0$ is interpreted by continuity.
\end{lemma}
\begin{proof}
The $j=k$ term of \eqref{res:classical} becomes
$\zeta(1+\varepsilon,a)$ after normalization. The gamma recurrence gives
\[
 \frac{k!t^{k+\varepsilon}\Gamma(-k-\varepsilon)}{(-t)^k}
 =-\frac{t^\varepsilon\Gamma(1-\varepsilon)}
       {\varepsilon\prod_{j=1}^k(1+\varepsilon/j)}.
\]
The Taylor series of
$\log[\Gamma(1-\varepsilon)/\prod_{j=1}^k(1+\varepsilon/j)]$ is
\[
 (\gamma-H_k)\varepsilon
  +\sum_{r=2}^{\infty}\frac{c_{k,r}}r\varepsilon^r.
\]
This proves \eqref{res:pair}; the remaining terms in the normally convergent
series are exactly \eqref{res:B}. Pairing the two poles first makes the
identity holomorphic at zero, so no value is assigned to either divergent
summand separately.
\end{proof}

\begin{theorem}[Uniform resonance, including growing integer order]
\label{res:main}
Fix $0<t_0<2\pi$, a compact set $A\subset(0,\infty)$, a compact set
$U\subset\C$, and integers $M,J\geq0$. Let $k\geq0$, $0<t\leq t_0$,
$a\in A$, $u\in U$, and $T$ be as in \eqref{res:T}. As $T\to\infty$,
uniformly over all these choices,
\begin{align}
 \mathcal R_k(u/T,t,a)
  ={}&T h(u)+\gamma_0(a)\notag\\
   &+\sum_{r=1}^{M}\frac{u^r}{T^r}
        \left\{\frac{(-1)^r\gamma_r(a)}{r!}
                  -e^{-u}b_{k,r+1}\right\}
        +O(T^{-M-1}),\label{res:mainformula}\\
 h(u)&=\frac{1-e^{-u}}u=\int_0^1e^{-uv}\dd v,\qquad h(0)=1.
 \label{res:h}
\end{align}
The error and its first $J$ derivatives in $u$ are $O(T^{-M-1})$. The same
claim holds after any fixed number of derivatives in $a$ on $A$.
Before absorbing the analytic tail into the error, it has the stronger bound
\begin{equation}\label{res:tailbound}
 B_k(u/T,t,a)=O\!\left(\frac{t}{k+1}\right)=O(e^{-T}).
\end{equation}
All constants may depend on the displayed compact sets and derivative
orders, but are independent of $k$ and $t$.
\end{theorem}

\begin{proof}
We give the uniformity argument, since it is what permits the integer order
to grow. Choose $0<\eta<1/2$. Since
$|c_{k,r}|\leq2\zeta(r)$ for every $k$ and $r\geq2$, the series
\[
 2\sum_{r=2}^{\infty}\frac{\zeta(r)}r\eta^r
\]
is a common upper bound for $|\log E_k(w)|$ on $|w|\leq\eta$.
Consequently the functions $E_k$ have Taylor remainders bounded uniformly
in $k$ on every smaller disk. The functions $G_a$ have the same property,
uniformly for $a$ in a compact complex neighborhood of $A$ contained in
$\Re a>0$.

It remains to control the infinite analytic tail uniformly in $k$. For
$\ell\geq1$,
\[
 \frac{k!}{(k+\ell)!}\leq\frac1{(k+1)\ell!}.
\]
Choose $t_0<R<2\pi$. Normal convergence of the regular part of the Lerch
series, uniformly for $|\varepsilon|\leq\eta$ and $a$ in the chosen compact
neighborhood, gives a finite common bound for
\[
 \sum_{\ell=1}^{\infty}
 \frac{t_0^{\ell-1}}{\ell!}
       |\zeta(1+\varepsilon-\ell,a)|.
\]
For an explicit justification of this parameter-uniform normal convergence,
use Hermite's Hurwitz integral \cite[25.11.29]{dlmfhurwitz}. Its integral
part is
\[
 i\int_0^\infty
 \frac{(a+iy)^{-s}-(a-iy)^{-s}}{e^{2\pi y}-1}\dd y.
\]
At $s=1+\varepsilon-\ell$, the sum of absolute majorants weighted by
$R^\ell/\ell!$ is bounded for $y\geq1$ by a fixed polynomial in $y$ times
$e^{-(2\pi-R)y}$. For $0\leq y\leq1$, the difference of the two powers
gains a factor $y$, which cancels the zero of the denominator; the remaining
factorial-weighted series is uniformly bounded. The two elementary terms
of Hermite's formula have normally convergent exponential majorants as well.
The compact neighborhood has $\Re a$ bounded away from zero, so the argument
also applies to complex $a$ and, by Cauchy's formula, to fixed derivatives.
This proves the required bound independently of $k$.

Thus $|B_k|\leq Ct/(k+1)$. The elementary inequality
$H_k-\gamma\leq\log(k+1)$ proves the second bound in
\eqref{res:tailbound}. For example, this harmonic inequality follows by
comparing the increasing sequence $H_k-\log(k+1)$ with its limit $\gamma$.

Now put $\varepsilon=u/T$ in \eqref{res:pair}. The first quotient is
\[
 T h(u)-e^{-u}\sum_{j=2}^{\infty}b_{k,j}(u/T)^{j-1}.
\]
Its apparent singularity at $u=0$ is removable. Uniform Taylor estimates
for this series and for $G_a(u/T)$ give \eqref{res:mainformula}; the
exponentially small tail is smaller than any specified power of $T^{-1}$.
To justify differentiated error estimates, first prove these bounds on
slightly larger compact sets in $u$ and $a$, then apply Cauchy's formula.
This avoids differentiating an unspecified asymptotic error.
\end{proof}

\begin{remark}[What the uniform statement covers]
For fixed $k$, the theorem describes $t\downarrow0$. It also holds when
$k\to\infty$ with $t$ fixed in $(0,2\pi)$, or along arbitrary simultaneous
sequences for which $T\to\infty$. The parameter $u$ remains in a fixed
compact set. The theorem does not assert the same error for unrestricted
$u\to\infty$, or uniformity as $a$ approaches zero.
\end{remark}

\subsection{Explicit corrections and the high-order coefficient limit}
The first coefficients are
\begin{align}
 b_{k,2}&=\tfrac12\bigl(\zeta(2)+H_k^{(2)}\bigr),\notag\\
 b_{k,3}&=\tfrac13\bigl(\zeta(3)-H_k^{(3)}\bigr),\label{res:firstb}\\
 b_{k,4}&=\tfrac14\bigl(\zeta(4)+H_k^{(4)}\bigr)
          +\tfrac18\bigl(\zeta(2)+H_k^{(2)}\bigr)^2.\notag
\end{align}
In particular,
\begin{align}
 \mathcal R_k(u/T,t,a)
  ={}&T h(u)-\psi(a)
 -\frac uT\left[\gamma_1(a)
        +\frac{e^{-u}}2\bigl(\zeta(2)+H_k^{(2)}\bigr)\right]\notag\\
 &+\frac{u^2}{T^2}\left[\frac{\gamma_2(a)}2
       -\frac{e^{-u}}3\bigl(\zeta(3)-H_k^{(3)}\bigr)\right]
       +O(T^{-3}).\label{res:explicit}
\end{align}

\begin{proposition}[The limiting coefficient function]\label{res:highorder}
Locally uniformly for $|w|<1$,
\begin{equation}\label{res:csc}
 E_k(w)\longrightarrow \Gamma(1-w)\Gamma(1+w)
             =\frac{\pi w}{\sin\pi w}.
\end{equation}
Hence $b_{k,2j+1}\to0$ and the even coefficients tend to those of
$\pi w/\sin\pi w$. The first nonconstant limits are
$b_{k,2}\to\pi^2/6$ and $b_{k,4}\to7\pi^4/360$.
\end{proposition}
\begin{proof}
For every $r\geq2$, $H_k^{(r)}\to\zeta(r)$. Dominated convergence in the
logarithmic series leaves $2\sum_{j\geq1}\zeta(2j)w^{2j}/(2j)$, which is
the logarithm of the gamma product in \eqref{res:csc}. The reflection
formula gives its sine expression. Local uniform convergence permits
coefficient extraction.
\end{proof}

This is a useful separation of roles. The generalized Stieltjes constants
carry the shift $a$; harmonic numbers and polygamma values carry the integer
order; the function $h$ carries the approach to resonance. The scale
$T=\log(1/t)+H_k-\gamma$ combines the two sources of logarithmic growth.

\subsection{Exact order derivatives and a Stieltjes finite part}
For later comparison with antiderivatives, set
\begin{equation}\label{res:gcoeff}
 \Gamma(1-w)=\sum_{j=0}^{\infty}g_jw^j,
 \quad g_0=1,\quad g_1=\gamma,\quad
 g_2=\frac{\gamma^2+\zeta(2)}2.
\end{equation}
For $L=\log(1/t)$ define the polynomial
\begin{equation}\label{res:P}
 P_n(L)=(-1)^{n+1}n!
       \sum_{j=0}^{n+1}g_j\frac{(-L)^{n+1-j}}{(n+1-j)!}.
\end{equation}
Its leading term is $L^{n+1}/(n+1)$. The first three are
\begin{align}
 P_0(L)&=L-\gamma,\notag\\
 P_1(L)&=\frac{L^2}{2}-\gamma L
                         +\frac{\gamma^2+\zeta(2)}2,\label{res:Plow}\\
 P_2(L)&=\frac{L^3}{3}-\gamma L^2
       +(\gamma^2+\zeta(2))L
       -\frac{\gamma^3+3\gamma\zeta(2)+2\zeta(3)}3.\notag
\end{align}

\begin{proposition}[Convergent finite-part identity]\label{res:finitepart}
For $a>0$, $n\geq0$, and $0<t<2\pi$,
\begin{align}
 e^{-at}\sum_{m=0}^{\infty}
       \frac{e^{-mt}\log^n(m+a)}{m+a}
  =P_n(\log(1/t))+\gamma_n(a)
       +(-1)^n\sum_{j=1}^{\infty}
          \zeta^{[n]}(1-j,a)\frac{(-t)^j}{j!}.
 \label{res:finitepartidentity}
\end{align}
The final series is convergent. The identity, including fixed $a$ derivatives,
holds locally uniformly. In particular the finite part at $t=0$ is
$\gamma_n(a)$, and the error after removing it is $O(t)$ on compact
positive $a$ intervals.
\end{proposition}
\begin{proof}
Put $k=0$ in the paired form of \eqref{res:classical}, expand at
$s=1+w$, and extract $(-1)^nn![w^n]$. The coefficient of
$[1-e^{-Lw}\Gamma(1-w)]/w$ is exactly $P_n(L)$.
The $j\geq1$ terms contain no pole near $w=0$, so Cauchy coefficient
extraction commutes with their normally convergent series.
\end{proof}

The proposition is an explicit generalized-shift version of the classical
integer-order derivative calculus, not a replacement for the new uniform
theorem. It supplies a useful numerical and conceptual bridge: the
subtraction is independent of $a$, so differentiation in $a$ removes it.
The unrenormalized $n=0$ Lerch value diverges at its endpoint, whereas its
finite part and parameter derivatives have canonical Hurwitz limits.

\subsection{Antiderivatives and exact logarithmic sum differences}
For $a,b>0$, subtracting \eqref{res:finitepartidentity} at the two shifts
cancels $P_n$. In particular,
\begin{align}
 \gamma_n(a)-\gamma_n(b)
 =\lim_{t\downarrow0}\sum_{m=0}^{\infty}
  \left[
  \frac{e^{-(m+a)t}\log^n(m+a)}{m+a}
 -\frac{e^{-(m+b)t}\log^n(m+b)}{m+b}\right].
 \label{res:shiftlimit}
\end{align}
For fixed $a,b$, the unweighted difference series also converges absolutely,
by the mean-value theorem and the bound
$O((1+\log^n m)/m^2)$ for its summand. The differently shifted exponential
weights have the same limit: differentiate their summand in the shift
variable and use $t e^{-(m+v)t}\leq1/[e(m+v)]$. This supplies the same
summable bound uniformly in $t>0$ on the interval between $a$ and $b$,
so dominated convergence applies. The integral in the shift variable is
\begin{equation}\label{res:primitive}
 \int_b^a\gamma_n(v)\dd v
  =\frac{(-1)^{n+1}}{n+1}
     \left[\zeta^{[n+1]}(0,a)-\zeta^{[n+1]}(0,b)\right].
\end{equation}
To see this without assigning values at a pole, integrate $G_v(w)$ first:
\begin{equation}\label{res:intG}
 \int_b^a G_v(w)\dd v
  =\frac{\zeta(w,b)-\zeta(w,a)-(a-b)}w.
\end{equation}
The numerator vanishes at zero because $\zeta(0,v)=1/2-v$; coefficient
comparison gives \eqref{res:primitive}. At $n=0$ this is
$-\log\Gamma(a)+\log\Gamma(b)$. At $n=1$ it is half the difference of
the second spectral derivatives at zero. These known primitive identities
are part of the manuscript's Hurwitz calculus \cite{repoIntegration}; the
new uniform resonance expansion supplies compatible limits when integer
order and argument are both changing.

\begin{remark}[A probability interpretation]
For $k=0$ and fixed $a$, the positive weights
$e^{-(m+a)t}/(m+a)$, normalized by their sum, define a probability law on
$\log(m+a)/\log(1/t)$. Its Laplace transforms tend to $h(u)$ for real $u$
in a neighborhood of zero, so the law converges to the uniform distribution
on $[0,1]$. This follows directly from Theorem~\ref{res:main} and the
standard uniqueness theorem for moment generating functions. The inverse
logarithmic corrections explain exactly where the Stieltjes constants and
gamma derivatives enter this limit. This interpretation is not needed in
any proof below.
\end{remark}
